% Suggested integration only; the structure owner chooses its location.
% The exact polynomial identities remain explicitly cited certificate inputs.
\begin{proposition}[Finite reconstruction of the parafermion factor]
\label{prop:pf13-finite-reconstruction}
Let $V,V'$ be simple vertex algebras of CFT type generated by a
conformal vector $L$ and a weight-three primary $a$. Set
$b=a_1a$ and $d=a_1b$, of weights four and five. Suppose the full
singular products $a_na,a_nb,a_nd$, $n\geq0$, agree as homogeneous
differential normal polynomials in $L,a,b,d$, with the same normal
ordering. Suppose also that their self-row is
\[
 a_5a=\frac c3\mathbf1,\quad a_4a=0,\quad a_3a=2L,\quad
 a_2a=TL,\quad a_1a=b,\quad a_0a=\frac12Tb-\frac1{12}T^3L,
\]
and that $a_nb=0$ for $n\geq4$, $a_3b\in\mathbb Ca$, and
$a_2b\in\mathbb CTa$. Then there is an isomorphism $V\cong V'$
identifying $L,a,b,d$. In particular $a_1d$ and both simple poles
must be specified polynomials; they cannot be left as independent
new generators.
\end{proposition}
\begin{proof}
We first recover the remaining generator OPEs without a circular
normal-product reduction. The Virasoro rows follow by commuting
Virasoro modes through $b=a_1a$ and $d=a_1b$:
\[
 L(m)b=(2m+1)a_{m+1}a,\qquad
 L(m)d=(2m+1)a_{m+1}b+a_1L(m)b\quad(m\geq1).
\]
These are linear in generators and derivatives under the displayed
hypotheses. Thus Virasoro creation operators can be moved to the
outside of any normal word using already known commutators.
Sesquilinearity removes derivatives. Since three non-Virasoro
factors have weight at least nine, the possible two-factor cores
through weight eight are exactly
\[
\begin{array}{c|c}
\text{weight bound}&\text{cores besides individual fields}\\\hline
6&:aa:\\
7&:aa:,\ :ab:\\
8&:aa:,\ :ab:,\ :bb:,\ :ad:.
\end{array}
\]
Reversing the order of a core uses skew symmetry and its already
known pair bracket. Commuting an external $a_n$ or $b_n$ through
outer Virasoro creation operators reduces their number; any resulting
negative mode is itself a normal product.

Write $A_r$ for the evaluation of $a_nP$, $n\geq0$, on every normal
polynomial $P$ of weight at most $r$, and $B_r$ for the corresponding
assertion with $b_n$. We use the noncommutative Wick identity
\[
 [x_\lambda:yz:]=:[x_\lambda y]z:+:y[x_\lambda z]:
              +\int_0^\lambda[[x_\lambda y]_\mu z]\,d\mu.
\]
The following order is acyclic. First $A_6$ is known: the only
two-factor core is $:aa:$, and $[a_\lambda a]$ is linear in the
vacuum and derivatives of $L,b$, whose brackets with $a$ are given.
The iterate identity
\begin{equation}
 b_ns=a_1(a_ns)-a_0(a_{n+1}s)
          +a_{n+1}(a_0s)-a_n(a_1s)
 \label{eq:pf13-b-iterate}
\end{equation}
with $s=b$ now determines $b_nb$, because each inner $a_jb$ has
weight at most six. Next $A_7$ follows: its new core $:ab:$ uses
only the just-computed $b,b$ row in the Wick integral.
Equation~\eqref{eq:pf13-b-iterate} with $s=d$ then determines
$b_nd$, since every inner $a_jd$ has weight at most seven.

At this point $B_6$ is known. For its core $:aa:$ the Wick integral
brackets a coefficient of $[b_\lambda a]$, of weight at most six,
with $a$, which is $A_6$ after skew symmetry. Single-field inputs
use the known $b,b$ and $b,d$ rows. We obtain $A_8$ next: its core
$:bb:$ uses $B_6$ on coefficients of $[a_\lambda b]$, and its core
$:ad:$ uses the known $L,d$ and $b,d$ rows. Finally $B_7$ follows
from $B_6$ on the coefficients of $[b_\lambda a]$ in the new core
$:ab:$. With $s=d$, the second iterate identity
\[
 d_ns=a_1(b_ns)-a_0(b_{n+1}s)
          +b_{n+1}(a_0s)-b_n(a_1s)
\]
computes $d_nd$: the first two terms use $A_8$, and the last two
use $B_7$. All generator singular products have therefore been
expressed by the same polynomials in the two algebras.

For the uniqueness step, take the vertex subalgebra $D\subset V\times V'$
generated by the diagonal pairs of the four named fields. Common
polynomial OPE closure and filtered Wick induction show that $D$ is
their differential normal span. Every generator has positive weight,
so $D_0=\mathbb C(\mathbf1,\mathbf1)$ and $D_n=0$ for $n<0$.
It follows that $D$ has a unique maximal proper graded ideal: every
proper graded ideal has zero degree-zero part, and so does their sum.
The coordinate projections are surjective because $L,a$ generate
each target. They preserve the conformal vector and hence the grading.
Their kernels are maximal proper graded ideals because $V,V'$ are
simple. The kernels coincide, which gives the asserted isomorphism.
No freeness or independence of normal monomials has been used.
\end{proof}

% Suggested application paragraph; bibliography keys must be selected by owner.
The exact parafermion records instantiate this proposition at
$c=8/5$ and $\lambda=14/473$. The upper-pole record supplies the
self-row and the remaining positive poles, the nonlinear record
supplies $a_0b$, and the truncation record supplies $a_1d$.
The algebraic simple-pole derivation supplies $a_0d$ from the
independently checked upper poles and skew symmetry. The
Dong--Lam--Yamada change-of-generators record matches the seven
weight-six decoupling coefficients with those of
$K(\mathfrak{sl}_2,13)$. The common first rows are the recursive
formulas (5.4), (5.7), and (A.2) of Linshaw. Simplicity of the actual
generated subalgebra follows from its PCT stability; its translation
operator agrees with the ambient one by the complementary positive
energy argument. Thus the proposition transports the checked
polynomial identities to an identification with the parafermion
algebra. The exact membership, PCT, translation, and polynomial
identities remain the explicit finite inputs to this identification.
